% Revised introduction: application, evidence trade-off, related approaches,
% specific gap, system overview, contributions, and organization.
\section{Introduction}\label{sec:intro}

\IEEEPARstart{U}{nmanned} aerial vehicles (UAVs) provide imagery for
inspection, traffic monitoring, and disaster assessment. Extracting useful
information from these images has motivated both human-assisted analysis
and automated object detection~\cite{ofli2016combining,pi2020convolutional}.
Visual question answering (VQA) offers another interface: users ask about
image content in natural language, such as the number of vehicles in a
scene. Remote sensing VQA has established this task and provided
benchmarks for evaluating it~\cite{lobry2020rsvqa}. With pretrained
vision--language models (VLMs)~\cite{wang2024qwen2vl}, an edge server can
process an aerial image together with a question. Applying this approach
over a UAV uplink, however, requires attention to both communication and
computation costs. Image transmission uses limited wireless resources,
while generating each answer also requires model inference. Reducing
the energy per question therefore calls for a design that considers both
costs.

The information required to answer a question need not include the full
image. For example, detected object classes and counts may support a
comparison between the numbers of cars and buses. Such structured
information has a small payload and can be processed by simple rules.
However, a missed object cannot generally be recovered from the detector
output alone. An image preserves visual details that may help answer
questions about objects omitted by the detector, but requires more
transmission and receiver computation. Moreover, the additional details
do not guarantee that a particular VLM will answer every question more
accurately. The choice depends on the question, the available detector
evidence, and the performance of the two processing paths over the link.
A useful selection method must identify when the smaller representation
is adequate and when image transmission is worth its cost.

Semantic and task-oriented communication provides the basis for making
such decisions according to the information needed by the receiver
task~\cite{gunduz2023beyond,strinati2021sixg,kountouris2021semantics}.
For wireless VQA, MU-DeepSC jointly designs the transceiver to convey
image and question features~\cite{xie2022mudeepscwcl}.
Goal-oriented VQA communication instead extracts and ranks bounding boxes
or scene-graph information according to the question~\cite{liu2024gosg}.
Query-oriented image coding also uses the query to select features for
image reconstruction~\cite{huang2026query}. Park and
Yoon~\cite{park2025transmit} compare visual representations across tasks
and select simpler or richer semantics according to task requirements.
Their task-level choice motivates examining which evidence and decoder
to use for each question within a VQA workload.
Du \etal~\cite{jiang2025lmmvn} use LLaVA for vehicle VQA and prioritize
image regions and patch transmission according to user interest.
That approach adapts the visual input to a multimodal model; our question
is when structured evidence can instead support symbolic decoding and
avoid VLM inference.
These studies establish task- and question-aware transmission as existing
design principles. In visual
reasoning, Prompt-RSVQA converts extracted visual context into text for a
language model~\cite{chappuis2022promptrsvqa}. Thus, structured visual
information is also an established alternative input to a reasoning
model.

A complementary line studies the energy implications of semantic
processing. Green semantic communication considers communication and
computation costs in aerial edge networks and probabilistic semantic
systems~\cite{zheng2024aerialee,dai2025pscom}. Input-aware VLM routing
selects between edge and cloud models according to the input and their
accuracy--cost trade-off~\cite{sabanovic2026inarvl}. Building on these
directions, we investigate a specific decision for structured aerial VQA:
whether to transmit detector information for symbolic decoding or an image
for VLM inference. The decision changes both the transmitted
representation and the computation needed to answer the question.
Its value must therefore be assessed through answer accuracy and the
combined communication and computation energy, rather than payload size
alone.

To study this decision, we propose an evidence-level semantic
communication system with a sender-side router and two receiver paths.
The token path transmits structured detector outputs and uses a
lightweight symbolic decoder. The image path transmits an image to a
frozen VLM. The system uses an aerial-domain detector and conventional
channel coding without jointly training the transceiver and the VLM.
A training-free rule selects the path from the question type. A learned
per-sample router further uses detector information to predict which path
will answer correctly, and an energy price controls the
accuracy--energy trade-off. We evaluate these choices on structured
questions about detected object categories. The aim is to identify when
evidence selection improves energy efficiency without reducing answer
accuracy, and to characterize the conditions under which the gain holds.
The main contributions are as follows:

\begin{itemize}
\item We formulate evidence selection as a joint choice of payload and
  receiver processing. The energy account includes transmission,
  onboard detection, and receiver inference, with UAV and edge costs
  distinguished. This formulation makes the cost of each answering path
  explicit and supports energy--accuracy comparisons.
\item We develop a training-free type rule and an MLP-based per-sample
  router. Conditional analysis explains when the type rule
  preserves accuracy while reducing energy. An energy-price extension
  provides a tunable operating point for the learned selector.
\item We evaluate the system on two aerial datasets, with AWGN,
  Rayleigh, and Rician channels on the primary dataset and Rician on
  the second. Evaluation combines logged-outcome replay, GPU power
  measurements, and transmission cost modeling. Under the reference
  configuration, the type rule reduces mean accounted energy per question by
  approximately $55\%$ on VisDrone relative to rate-adaptive image
  transmission while improving accuracy. On DroneVehicle, a
  validation-selected learned router achieves $74.60\%$ accuracy
  at $4.47$ J per query, compared with $72.93\%$ at $11.21$ J
  for the type rule. We retain simple baselines and matched-receiver
  comparisons to identify where learned selection helps and where
  its additional complexity offers little benefit. These results
  concern system energy, not a measured UAV battery gain.
\end{itemize}

The remainder of this paper is organized as follows.
Section~\ref{sec:related} reviews related work.
Section~\ref{sec:system} presents the system model and energy accounting.
Section~\ref{sec:routing} develops the selectors and their analytical
support. Section~\ref{sec:exp} reports the experiments and discusses
their limitations. Section~\ref{sec:conc} concludes the paper.
